\section{Discussion and Limitations}
\label{sec:discussion}

Our results separate three requirements for reasoning through a visual intermediate state: the state must be \emph{correct in its relations}, though not necessarily precise; the model must \emph{read} it, which it does well on viewpoint and stability and less well on rearrangement; and the model must \emph{rely} on it rather than on the original image. Current VLMs meet the second requirement, but not the first or the third. Text reasoning is not always worse: when a question can be rephrased in terms of what is already visible, it does as well or better, at a much higher token cost. A system that routes between the two would need to recognise which case it is in. Beyond routing, two directions follow: training signals for relational accuracy of the imagined state, and mechanisms that let the imagined state override the observation it was derived from.

\paragraph{Limitations.} Our scenes are synthetic and simple, which is what makes the oracle exact; real scenes would add perception difficulty on top. The physical aftermath is a plausible depiction, not a simulation, and for stable stacks the imagined image equals the original. Given only the aftermath, models answer with a fixed prior, so we do not interpret imagined-only accuracy on that task. We evaluate two open 7--8B models with 4-bit weights; larger models may rely on the imagined state differently.
